\begin{alist}
    \item Στο τρίγωνο $\mathrm{B}\mathrm{E}\Gamma$ το τμήμα $\mathrm{E}\mathrm{M}$ είναι ύψος και διάμεσος, οπότε το τρίγωνο $\mathrm{B}\mathrm{E}\Gamma$ είναι ισοσκελές.

    \item Αφού $\Theta, \mathrm{Z}$ είναι οι προβολές του $\mathrm{E}$ στις $\mathrm{A}\mathrm{B}, \mathrm{A}\Gamma$ αντίστοιχα τότε τα $\mathrm{E}\Theta, \mathrm{E}\mathrm{Z}$ θα είναι κάθετα τμήματα στις $\mathrm{A}\mathrm{B}, \mathrm{A}\Gamma$ οπότε $\mathrm{E}\hat{\Theta}\mathrm{B} = \mathrm{E}\hat{\mathrm{Z}}\Gamma = 90^\circ$. \\
    Τα τρίγωνα $\mathrm{E}\Theta\mathrm{B}$ και $\mathrm{E}\mathrm{Z}\Gamma$ έχουν:
    \begin{itemize}
        \item $\mathrm{E}\hat{\Theta}\mathrm{B} = \mathrm{E}\hat{\mathrm{Z}}\Gamma = 90^\circ$
        \item $\mathrm{E}\mathrm{B} = \mathrm{E}\Gamma$, ως πλευρές του ισοσκελούς τριγώνου $\mathrm{E}\mathrm{B}\Gamma$ από α) ερώτημα.
        \item $\mathrm{E}\Theta = \mathrm{E}\mathrm{Z}$, γιατί το $\mathrm{E}$ είναι σημείο της διχοτόμου $\mathrm{A}\mathrm{E}$ οπότε ισαπέχει από τις πλευρές της γωνίας $\mathrm{B}\hat{\mathrm{A}}\Gamma$.
    \end{itemize}
    Άρα, τα τρίγωνα $\mathrm{E}\Theta\mathrm{B}$ και $\mathrm{E}\mathrm{Z}\Gamma$ είναι ίσα ως ορθογώνια που έχουν δυο ομόλογες πλευρές τους (υποτείνουσα - κάθετη πλευρά) ίσες μία προς μία.

    \item Επειδή τα τρίγωνα $\mathrm{E}\Theta\mathrm{B}$ και $\mathrm{E}\mathrm{Z}\Gamma$ είναι ίσα θα ισχύει και $\Theta\hat{\mathrm{B}}\mathrm{E} = \mathrm{Z}\hat{\Gamma}\mathrm{E}$ ή διαφορετικά $\Theta\hat{\mathrm{B}}\mathrm{E} = \mathrm{A}\hat{\Gamma}\mathrm{E} \quad (1)$. \\
    Ισχύει ότι $\Theta\hat{\mathrm{B}}\mathrm{E} + \mathrm{A}\hat{\mathrm{B}}\mathrm{E} = 180^\circ$ (ως παραπληρωματικές) και αφού $\Theta\hat{\mathrm{B}}\mathrm{E} = \mathrm{A}\hat{\Gamma}\mathrm{E}$ τότε θα είναι $\mathrm{A}\hat{\Gamma}\mathrm{E} + \mathrm{A}\hat{\mathrm{B}}\mathrm{E} = 180^\circ$.
\end{alist}
